\chapter{First Setup}\label{ch:setup}

Do the procedures in this chapter in the given sequence when you start a new
machine. You can do them in the console of the board, or in
\menu{Plotter settings} in the desktop app (section~\ref{sec:app-settings}).

\section{Open the Console}\label{sec:console}

The console is a text interface for setup and for diagnosis. It is available
on two links:

\begin{description}[style=nextline,leftmargin=1.2em,font=\sffamily\bfseries\color{navy}]
  \item[USB] Open the USB COM port in a serial terminal. Type \cmd{menu} on a
        line by itself and press Enter. Type \cmd{exit} to go back to HP-GL
        mode. While the console is open, the USB port does not receive HP-GL.
  \item[TCP port 9300] Connect a terminal (for example PuTTY in ``raw'' mode,
        or \code{ncat}) to port 9300 of the board. The console accepts one
        client at a time. The desktop app uses this port when it is connected
        over Ethernet.
\end{description}

Each command ends with a line that starts with \code{ok} or \code{error:}.
Appendix~\ref{app:console} lists all commands. These are the commands that
you use most:

\begin{lstlisting}[language=console]
status                 # state, position, pen, counters, and links
config                 # all settings and their values
set <key> <value>      # change a setting and save it
enable / disable       # supply or remove motor power
goto <x_mm> <y_mm>     # pen-up move to a position
pen up / pen down      # test the pen lift
origin                 # the present position becomes plotter 0,0
clear                  # clear a latched fault
\end{lstlisting}

\subsection{When a Setting Takes Effect}

Each setting has a rule for when you can change it. The \cmd{config} command
shows the rule at the end of the line. The desktop app shows it as a label
next to the field.

\begin{table}[htbp]
\centering
\caption{When a setting can change}\label{tab:when}
\begin{tabularx}{\linewidth}{@{}P{0.17\linewidth}P{0.2\linewidth}L@{}}
\toprule
\textbf{Mark in \cmd{config}} & \textbf{Label in the app} & \textbf{Rule} \\
\midrule
(none) & (none) & Live. The change takes effect immediately, also during a job. \\
\code{[idle]} & between jobs & The board refuses the change while a job plots. The motors can stay enabled. \\
\code{[disabled]} & motors disabled & The board refuses the change while the motors are enabled. Use \cmd{disable} first. The app does this for you. \\
\code{[restart]} & after restart & The board saves the change. It takes effect after \cmd{restart}. \\
\bottomrule
\end{tabularx}
\end{table}

\section{Set the Kinematics}

Set the kinematics first. It is the first setting in \cmd{config} and in the
app, because it selects the mechanics settings that apply.

\begin{itemize}
  \item \key{kinematics}~0, Cartesian (default): one motor for each axis.
        \conn{M0} moves X and \conn{M1} moves Y.
  \item \key{kinematics}~1, CoreXY or H-bot: two belts move the head together
        (section~\ref{sec:corexy}).
\end{itemize}

\begin{procedure}{Set the kinematics}
\begin{enumerate}
  \item Type \cmd{disable}.
  \item Type \cmd{set kinematics 0} for a Cartesian machine, or
        \cmd{set kinematics 1} for a CoreXY or H-bot machine. The present
        position becomes the origin.
\end{enumerate}
\end{procedure}

\section{Set the Mechanics}

The board must know how far the head moves for one revolution of a motor.

\begin{itemize}
  \item \key{steps_rev_x}, \key{steps_rev_y}: the steps per revolution of
        \conn{M0} and \conn{M1}. Use the input resolution that you set in
        ClearPath MSP. Default 800.
  \item Cartesian: \key{mm_rev_x}, \key{mm_rev_y}: the travel of the axis for
        one revolution of the motor, in mm. Default 5.0. For a lead screw, it
        is the lead of the screw. For a belt, it is the number of pulley teeth
        multiplied by the belt pitch.
  \item CoreXY: \key{pulley_dia_mm}: the pitch diameter of the pulleys, in
        mm, for both motors. Default \SI{12.7324}{mm} (a 20-tooth GT2 pulley).
        The pitch diameter is the number of teeth multiplied by the belt
        pitch, divided by $\pi$.
\end{itemize}

\begin{procedure}{Set the mechanics of a CoreXY machine}
\begin{enumerate}
  \item Type \cmd{disable}.
  \item Type \cmd{set steps_rev_x 800} and \cmd{set steps_rev_y 800}. Use the
        value from ClearPath MSP. The two values must be the same.
  \item Type \cmd{set pulley_dia_mm 12.7324}. Use the pitch diameter of your
        pulleys.
  \item Type \cmd{config}, and make sure that the values are correct.
\end{enumerate}
\end{procedure}

\begin{procedure}{Set the steps and the travel per revolution (Cartesian)}
\begin{enumerate}
  \item Type \cmd{disable}.
  \item Type \cmd{set steps_rev_x 800}. Use the value from ClearPath MSP.
  \item Type \cmd{set mm_rev_x 40}. Use the travel of your X axis.
  \item Do steps 2 and 3 again for Y, with \key{steps_rev_y} and \key{mm_rev_y}.
  \item Type \cmd{config}, and make sure that the four values are correct.
\end{enumerate}
\end{procedure}

\begin{note}
\notelabel On a Cartesian machine, use the same steps per millimeter on X and
Y. If the two axes are different, the speeds use the X scale, and arcs are
drawn as short straight lines. On a CoreXY machine, the scales are always the
same.
\end{note}

\section{Set the Plot Area and the Origin}

The \term{plot area} is \key{width_mm} by \key{height_mm} (default
400 by 300). It is the \term{hard clip}: the board does not draw outside it.
Plotter 0,0 is the lower-left corner of the plot area. The position of the
pen at power-up is 0,0, until you set a different origin.

\begin{warning}
The board does not know where the end stops of your machine are. Set the plot
area to the travel that is really available from the origin. If the area is
too large, the carriage can hit the end of its travel.
\end{warning}

\begin{procedure}{Set the plot area and the origin}
\begin{enumerate}
  \item Measure the travel of each axis from the corner that you want as the
        origin.
  \item Type \cmd{set width_mm 600}. Use your X travel.
  \item Type \cmd{set height_mm 300}. Use your Y travel.
  \item Move the pen to the lower-left corner of the plot area. Use
        \cmd{goto} or the \menu{Jog} tab of the app.
  \item Type \cmd{origin}. The present position becomes plotter 0,0.
\end{enumerate}
\end{procedure}

\begin{note}
\notelabel The board does not keep the origin when it restarts. After power-up,
the origin is where the pen is. Move the pen to the corner before you supply
power, or set the origin again.
\end{note}

\section{Make Sure That the Axes Move Correctly}

\begin{warning}
Keep your hands out of the work area. The axes move when you send the next
commands.
\end{warning}

\begin{procedure}{Do the first moves}
\begin{enumerate}
  \item Type \cmd{set draw_mm_s 20} and \cmd{set travel_mm_s 20} to get a low speed.
  \item Type \cmd{enable}. Make sure that the reply is \code{ok}.
  \item Type \cmd{goto 50 0}. The pen must move \SI{50}{mm} in the positive
        X direction (to the right on a standard plotter).
  \item Measure the distance. If it is not \SI{50}{mm}, examine
        \key{steps_rev_x} and \key{mm_rev_x}.
  \item If the axis moved in the wrong direction, type \cmd{goto 0 0},
        \cmd{disable}, then \cmd{set dir_x -1}.
  \item Do steps 3 to 5 again for Y with \cmd{goto 0 50} and \key{dir_y}.
  \item Type \cmd{goto 0 0} to go back to the origin.
\end{enumerate}
\end{procedure}

\subsection{CoreXY and H-Bot Machines}\label{sec:corexy-check}

On a CoreXY or an H-bot machine (section~\ref{sec:corexy}), set the
kinematics first. Then find the correct directions with this procedure.

\begin{warning}
With the wrong kinematics or directions, the head can move in a direction
that you do not expect. Use a low speed and short moves, and keep the head
away from the ends of its travel.
\end{warning}

\begin{procedure}{Set the kinematics and the directions of a CoreXY machine}
\begin{enumerate}
  \item Type \cmd{disable}.
  \item Make sure that \key{kinematics} is 1 and that the steps per
        revolution and \key{pulley_dia_mm} are set.
  \item Type \cmd{set draw_mm_s 20} and \cmd{set travel_mm_s 20}.
  \item Type \cmd{enable}, then \cmd{goto 20 0}.
  \item Look at the movement of the head, and find it in
        table~\ref{tab:corexy-fix}.
  \item Type \cmd{disable}, then do the correction in the table.
  \item Type \cmd{origin}, and do steps 4 to 6 again until the head moves
        \SI{20}{mm} to the right.
  \item Type \cmd{goto 20 20}. The head must move \SI{20}{mm} up (positive Y).
        If it moves down, type \cmd{disable}, set \key{dir_y} to the other
        value, type \cmd{origin}, and do this step again from \cmd{goto 20 0}.
\end{enumerate}
\end{procedure}

\begin{table}[htbp]
\centering
\caption{Corrections for a CoreXY machine after \cmd{goto 20 0}}\label{tab:corexy-fix}
\begin{tabularx}{\linewidth}{@{}P{0.38\linewidth}L@{}}
\toprule
\textbf{The head moves} & \textbf{Correction} \\
\midrule
Right (positive X) & Correct. Go to the last step. \\
Left (negative X) & Set \key{dir_x} to the other value (1 or -1). \\
Up or down (Y) & Set \key{swap_xy} to the other value (0 or 1). \\
Diagonally & \key{kinematics} is not 1, or the belts are not a CoreXY or H-bot layout. \\
\bottomrule
\end{tabularx}
\end{table}

\section{Set the Pen Lift}

\begin{procedure}{Set and test the pen lift}
\begin{enumerate}
  \item Type \cmd{disable}.
  \item Type \cmd{set pen_mode 0} for a digital output, or \cmd{set pen_mode 1}
        for a servo on \conn{M2}.
  \item For a digital output, set \key{pen_io} (0 to 5) and \key{pen_active_high}.
  \item For a servo on \conn{M2}, set \key{z_up_steps} and \key{z_down_steps}.
  \item Type \cmd{enable}.
  \item Type \cmd{pen down}. Make sure that the pen goes down.
  \item Type \cmd{pen up}. Make sure that the pen goes up.
  \item If the pen moves before it is fully down or up, increase
        \key{pen_down_ms} or \key{pen_up_ms}.
\end{enumerate}
\end{procedure}

\section{Set the Speeds}

Table~\ref{tab:speeds} shows the motion settings. All of them are live, so
you can change them during a test plot and see the result.

\begin{table}[htbp]
\centering
\caption{Motion settings}\label{tab:speeds}
\begin{tabularx}{\linewidth}{@{}P{0.21\linewidth}P{0.15\linewidth}L@{}}
\toprule
\textbf{Setting} & \textbf{Default} & \textbf{Function} \\
\midrule
\key{draw_mm_s} & 100 mm/s & Speed with the pen down. \\
\key{travel_mm_s} & 200 mm/s & Speed with the pen up. \\
\key{max_mm_s} & 250 mm/s & Highest pen-down speed that a job can ask for with \cmd{VS}. \\
\key{job_speed} & 1 & 1: the speed in the job (\cmd{VS}) sets the pen-down speed. 0: always use \key{draw_mm_s}. \\
\key{accel_mm_s2} & 2000 mm/s\textsuperscript{2} & Acceleration along the path. It also limits the speed on curves. \\
\key{corner_mm_s} & 20 mm/s & Largest change of speed of one motor at a corner. A larger value gives faster corners and more shock to the machine. \\
\key{motor_max_mm_s} & 0 (none) & Highest speed of a motor, as belt travel. For CoreXY, where a motor can turn 1.41 times faster than the head moves. \\
\bottomrule
\end{tabularx}
\end{table}

Start with low values. Increase them until the machine is fast and the lines
are still clean. A heavy carriage or a long belt needs a lower acceleration.

\section{Set the Network}

The board uses DHCP by default. If it gets no address, it uses
192.168.0.110. For a fixed address, do this procedure.

\begin{procedure}{Set a static IP address}
\begin{enumerate}
  \item Type \cmd{set net_mode 1}.
  \item Type \cmd{set ip 10.168.41.5}. Use your address.
  \item Type \cmd{set netmask 255.255.255.0}.
  \item Type \cmd{set gateway 10.168.41.1}. If no other network must reach the
        board, you can keep 0.0.0.0.
  \item Type \cmd{restart}. The board restarts with the new address.
  \item From the computer, type \cmd{ping 10.168.41.5} to make sure that the
        board replies.
\end{enumerate}
\end{procedure}

To go back to DHCP, type \cmd{set net_mode 0} and \cmd{restart}.

\section{Do a Test Plot}

Send a small square to the board. On Windows, save this text in a file with
the name \file{square.plt}:

\begin{lstlisting}[language=hpgl]
IN;SP1;PU400,400;PD2400,400,2400,2400,400,2400,400,400;PU0,0;
\end{lstlisting}

This draws a \SI{50}{mm} square that starts \SI{10}{mm} from the origin.
Send it to the USB port (for example COM3) with this command in
\file{cmd.exe}:

\begin{lstlisting}[language=console]
copy /b square.plt COM3
\end{lstlisting}

The motors enable automatically, the square plots, and the pen goes back to
0,0. Measure the square. Each side must be \SI{50}{mm}.

\begin{note}
\notelabel The board keeps a job from the USB port until the port is quiet
for 2~s, and then plots it in a better order (section~\ref{sec:optimizer}).
Thus the plot starts approximately 2~s after the file is sent.
\end{note}
